\subsection{Carry the historical terms through a complete command}
\label{sec:selected-term-execution}

The historical checks in Section~\ref{sec:nested-term-invariants} answer whether
the supplied terms and assessments fit together. They do not, by themselves,
justify changing the claim. An amendment also needs a valid current state,
an identified customer and prior term, an admissible effective date, and a
complete description of the writes it will perform. A collection must change
the remaining receivable without rewriting its earlier assessments. At a
period boundary, closing entries and the next assessment window must agree.
These are connected requirements: accepting each snapshot in isolation would
leave the correspondence between the event and the financial history unchecked.

The selected continuation described here joins those checks for the retained
synthetic claim. It covers prospective due-date amendments, collections,
assessments and period carry. The accounting premise is still deliberately
narrow: an amendment changes the supplied terms without recognising a
modification gain or loss, interest or impairment. Actual contractual
applicability and the appropriate accounting treatment require their own
evidence. Within this premise, the implementation question is precise: can
explicit declarations produce the complete retained requests and results while
the native financial builders are disabled?

\paragraph{Establish which supplied state the declaration accepts.}
The selected state contract combines structural checks with the six historical
and financial checks already described. Structural checks require the exact
state fields and profile, a typed book, the declared term-claim and clock
schemas, the operating account roles, a positive cost leaf with the expected
units, and the stated migration fields and premise. The account selection is
explicit. A different cost pointer can be valid under the native validator
and still fall outside this selected contract; refusal then means that the
selection does not cover that configuration.

Existing admission stages evaluate most of these conditions on a normalized
view in which the original \code{book} field is named \code{opening}. The
original envelope is checked first, and normalization preserves every value
and its insertion order. The reused stage engine's internal success label does
not turn a mid-period snapshot into a new-period opening. The enclosing result
reports satisfaction of the selected snapshot contract and retains false
opening, source and accounting admission flags.

Two additional conditions are explicit. The migration source-book identifier
must have the syntax of a SHA-256 digest, and every term premise must contain
nonblank text. The distinction between text and nonblank text matters: a
space-only premise can satisfy the generic text type while failing the native
premise requirement. Checking a digest's syntax establishes neither the
existence nor the truth of the referenced source. Similarly, complete supplied
claim and nested-row populations remain a stated input scope. A partial
population leaves the command unresolved before construction.

\paragraph{Resolve the event against the same history.}
For an amendment, eligibility first checks the event's identity and date, the
complete typed inputs, and the supplied approval. Unknown inputs produce an
unresolved result. An explicitly unselected approval produces an unsupported
result. These outcomes differ from invalid state, stale selection, duplicate
event or contradictory dates.

The event must name an existing claim belonging to the stated customer. Its
prior-term input must equal the claim's current term. The new term identifier
must be nonblank, unused and different from the sentinel \code{none}; the
premise must also be nonblank. The new due date cannot precede the event's
effective date. Finally, the effective date must be strictly later than that
of the selected prior term. Equal dates would make the prospective sequence
ambiguous and are refused, even when the supplied snapshot is otherwise valid.

A nested lookup is permitted only after the declaration establishes the
presence of the parent claim and the selected child term. This is an ordering
requirement on the checks, not an inference from a successful dictionary
lookup. The following binding fragment names the prior term's effective date;
the two referenced stages must appear earlier and establish the corresponding
memberships. The surrounding rule declares its inputs and comparison.

\begin{verbatim}
binding "previous";
  kind nested_map_field;
  pointer "/claims";
  key_input "claim_id";
  parent_after_stage "claim_exists";
  child_pointer "/terms";
  child_key_input "prior_term";
  field "effective";
  after_stage "prior_exists";
end;
\end{verbatim}

The selected eligibility engine checks those dependencies when compiling the
rule. It then resolves the bound values against the actual book and event.
The preserved guarded-predicate evaluator owns the comparisons. Adding nested
membership and field lookup does not introduce a second arithmetic or date
language.

\paragraph{Construct the complete nested change.}
An eligible amendment constructs three ordered memory operations: insert the
new term record, replace the claim's current-term identifier, and replace its
current due date. The inserted record states due date, effective date, event
identifier, predecessor and premise. Each value comes from an explicitly typed
parameter. The parameterized constructor builds the full method and typed
transition request; it does not import a stored native method as the answer
for the current event.

An assessment has a different construction. For every supplied claim, it
inserts a complete assessment record and replaces five summary fields:
latest assessment, assessed flag, assessment date, assessed remaining amount
and overdue amount. The record retains the applicable term identifier, so
later amendments cannot change the basis of this assessment. A final operation
advances the assessment cursor to the end of the checked window. In the
retained one-claim case this produces seven writes. Two claims produce thirteen
writes, with a single sequence of operation identifiers across both claims and
the suffix.

The previous ordered expander admitted only scalar fields in each row and
scalar replacements. It could not represent an insertion into the nested
assessment map. Its successor therefore checks a declared nested row schema
and recursively checks complete write templates. A record path must name a
declared field; a dynamic map key must be typed text. Insertions target map
members and supply the entire declared record. Replacements target scalar
leaves with the correct type. Both conditional branches are checked, including
the branch not taken by the current claim.

Validation precedes expansion even when the claim map is empty. An empty map
then emits only the cursor suffix, while an invalid nested record template
still fails. The existing 256-write limit is retained: a construction that
would emit 257 writes refuses as a whole. This limit is an implementation
boundary, not a conclusion that a larger company's history is economically
impossible. Claim iteration order is explicit and is preserved in the writes;
it does not determine which term was effective on an earlier date.

\paragraph{Check the proposed result before continuing.}
The command composition checks the selected snapshot before eligibility and
again after the typed transition. A well-typed but incorrect predecessor or
overdue write can pass construction and still fail the resulting historical
requirements. In that case the result retains diagnostic transition evidence
but does not return an accepted top-level after-state. The caller cannot
continue the sequence from that rejected proposal. The successful path remains
a provisional calculation; it performs no durable commit and does not establish
source or accounting admission.

Period carry follows the same discipline. Its report request must bind to the
exact closing book. The selected policy preserves all typed memory leaves,
including immutable term and assessment histories, and resets only the
assessment cursor to the next period's start. The preserved report and carry
engines produce the closing and opening books, which are checked against the
selected snapshot contract. Completeness of the supplied population is carried
forward only from the accepted pre-state through this declared transition;
it is not manufactured for an unrelated supplied book.

Chronology is visible in the complete comparison. Each company's retained
sequence contains nine events, with carry between the seventh and eighth
events. The continuation uses the computed output of each command as the next
input, including carry. Appending carry after all nine events, or resetting
the input to a stored snapshot before each event, would not verify this chain.
The final state agrees only after the commands are placed at their actual
positions across the two fiscal periods.

\paragraph{Make the selected behavior directly authorable.}
The state-contract source spells out structural stages, nested projections,
queries, captures, template substitutions, guarded requirements and collection
constraints. Event-rule source separately spells out eligibility bindings,
typed parameters, method conditions, monetary effects and nested writes.
Carry policy has its own short declaration. These source forms contain the
algorithm-bearing declarations rather than an opaque import of a completed
financial algorithm. JSON tokens remain available for explicit types, contexts,
schema expectations and scalar data.

The frontend preserves authored declaration order and exact Unicode strings.
Only line-feed characters split source lines; carriage-return/line-feed input
is also accepted without mistaking a Unicode line separator inside a string
for a new declaration. Locations record physical line numbers and UTF-8 byte
slices. Lexical errors identify tokens or declarations; failures reported by
an opaque contextual validator identify the responsible source component,
without pretending to identify a more precise field.

Parsing, eligibility compilation and contextual construction remain distinct
states. A parsed construction is not reported as contextually compiled until
its event and book have been supplied. Execution also requires the independently
supplied selection of the combined contract and rules. Changing a source file
does not authorize its changed declaration: the old selection refuses it.
Source diagnostics are retained alongside the complete typed execution result,
so adding location information does not obscure the comparison with that result.

The finite evidence establishes this selected continuation. Structural checks
and direct state-contract source each reproduce 74 complete comparisons. The
nested constructor reproduces ten retained assessment requests and transitions.
Full command composition and its source counterpart each reproduce two
continuous traces, comprising eighteen events and two carries. Malformed
sources, incorrect bindings, partial populations, stale selections, nested
write errors and the equal-effective-date boundary have separate controls.
Native financial callbacks are disabled during candidate execution. Shared
parser and typed-engine defects remain possible; this is bounded conformance
evidence, not independent acceptance or whole-company coverage.

The complete sources, expected declarations, failures and results are retained
under \path{data/accounting-master/term-structural-admission-20261026/},
\path{data/accounting-master/nested-contract-source-20261026/},
\path{data/accounting-master/nested-assessment-construction-20261026/},
\path{data/accounting-master/term-full-event-binding-20261026/} and
\path{data/accounting-master/term-event-source-20261026/}.
